\section{Historical projection backtest}

The retrospective analysis uses archived WPP revisions whose target years have
since elapsed and are covered by a later WPP estimate series. Each vintage is
evaluated only over horizons that had resolved by the comparison date. The
current implementation covers the 1992, 1994, 1996, 2000, 2002, 2004, 2006,
and 2008 revisions and reports results at the world and stable regional levels,
avoiding country comparisons where changing borders would create false errors.
The omitted 1998 release and later revisions will be documented explicitly
when the historical archive is extended; the present eight-vintage set is the
implemented and reproducible scope rather than a claim to cover every release
in the period \citep{unwpparchives}.

The comparison series is WPP 2024. It is an updated demographic estimate, not
ground truth. This distinction is especially important for weak-data countries
and much less consequential for world totals. The low and high projection
variants are fertility scenarios rather than probability intervals; whether a
later estimate falls between them is therefore described as inclusion in the
printed range, not probabilistic calibration.
The current calculation conditions on the WPP 2024 estimates as fixed
comparison values and does not propagate uncertainty in those estimates.

For population, the analysis records signed and absolute relative error by
vintage, target year, horizon, and region. Fertility and life expectancy are
evaluated separately to identify which projected components moved in the same
direction as population error. Overlapping target years and successive
vintages are not independent observations. The final paper will therefore
present vintage-level distributions and robustness analyses rather than treat
every vintage-year pair as an independent replication.

Across the scored world-level vintage-year pairs, the mean signed population
error is \BacktestWorldMeanSignedError, where a negative value means that the
projection was below the later WPP 2024 estimate. The 1992 revision projected
too high; every included revision from 1996 onward projected too low. The
forecasts overlap in target years and are not independent tests of an unchanged
model.

\begin{figure}[htbp]
  \centering
  \includegraphics[width=\textwidth]{backtest-world-population.png}
  \caption{Medium-variant world population projections from eight archived WPP
  revisions, shown only for target years available in the later WPP 2024
  estimate series. The right panel reports signed relative error; negative
  values are underprojections.}
  \label{fig:world-backtest}
\end{figure}

The component backtest points in different directions across regions. Mean
African fertility error is \BacktestAfricaTfrError, while mean Eastern Asian
fertility error is \BacktestEastAsiaTfrError. Life expectancy is lower than the
later comparison series by \BacktestLifeExpectancyError{} on average. These
figures identify where the archived trajectories missed; they do not establish
a causal mechanism. In a separate range check, \BacktestRangeInclusion{} of the
world-level observations fall inside the printed low-to-high fertility
variants. That quantity is not a probability-interval coverage rate.

\begin{figure}[htbp]
  \centering
  \includegraphics[width=0.95\textwidth]{backtest-fertility.png}
  \caption{Archived total-fertility projections for Africa and Eastern Asia
  compared with WPP 2024 estimates. The opposite-signed misses motivate a
  mechanism-explicit comparison, but do not themselves attribute cause.}
  \label{fig:fertility-backtest}
\end{figure}

A formal component decomposition and extension to later gradeable revisions
remain required before making stronger causal claims.
